\section{养殖业：生猪、肉鸡与其他养殖}
\label{sec:breeding}

养殖业是农业上市公司中市值最大的板块：\num{22} 家公司、流通市值合计约 \num{3603} 亿元，
占全部农业上市公司合计流通市值的 \hl{38.3\%}。
2025 年 \num{9} 家亏损、\num{6} 家资产负债率超过 \num{65}\%，
ROE 中位数 \num{2.64}\%。\hl{这是整个报告中最需要“看结构而不是看平均”的板块}：
平均值掩盖了牧原 \num{20.6}\% 与天邦 \num{-13.1} 亿元之间的巨大差异。

\begingroup\scriptsize\setlength{\tabcolsep}{2pt}
\begin{longtable}{@{}llrrrrrrr@{}}
\caption{农业上市公司分行业明细（生猪养殖、肉鸡养殖、其他养殖；财务为 2025 年年报，公告计数窗口 2023-09---2026-09）}\label{tab:comp-breeding}\\
\toprule
代码 & 公司简称 & 市值（亿元） & 营收（亿元） & 净利（亿元） & CAGR（\%） & ROE（\%） & 再融资公告 & 重组公告 \\
\midrule
\endfirsthead
\multicolumn{9}{l}{\small（续）}\\
\toprule
代码 & 公司简称 & 市值（亿元） & 营收（亿元） & 净利（亿元） & CAGR（\%） & ROE（\%） & 再融资公告 & 重组公告 \\
\midrule
\endhead
\bottomrule
\endlastfoot
002714 & 牧原股份 & 1,327.9 & 1,441.4 & 154.87 & 14.0 & 20.6 & 23 & 0 \\
300498 & 温氏股份 & 862.7 & 1,038.6 & 52.66 & 7.5 & 12.6 & 58 & 3 \\
002157 & 正邦科技 & 227.4 & 147.9 & -5.46 & 45.5 & -4.9 & 12 & 0 \\
002299 & 圣农发展 & 191.3 & 200.9 & 13.80 & 4.3 & 12.6 & 0 & 3 \\
605296 & 神农集团 & 156.9 & 53.5 & 3.39 & 17.3 & 6.9 & 27 & 1 \\
300761 & 立华股份 & 132.1 & 187.3 & 5.78 & 10.4 & 6.3 & 3 & 0 \\
000048 & 京基智农 & 98.8 & 48.7 & 1.53 & -37.4 & 3.6 & 9 & 1 \\
002458 & 益生股份 & 78.2 & 29.5 & 1.65 & -4.3 & 3.8 & 10 & 0 \\
603477 & 巨星农牧 & 70.8 & 80.0 & 0.30 & 40.7 & 0.9 & 75 & 0 \\
000735 & 罗牛山 & 64.9 & 33.8 & 1.54 & -10.8 & 3.6 & 1 & 4 \\
600975 & 新五丰 & 63.5 & 72.6 & -8.76 & 13.5 & -33.9 & 4 & 3 \\
002321 & 华英农业 & 48.2 & 54.5 & -0.60 & 21.3 & -6.4 & 1 & 0 \\
002746 & 仙坛股份 & 43.9 & 57.7 & 2.43 & -0.1 & 5.3 & 0 & 2 \\
301116 & 益客食品 & 40.9 & 189.6 & -2.80 & -6.9 & -17.2 & 48 & 0 \\
002124 & 天邦食品 & 40.8 & 86.9 & -13.09 & -7.8 & -43.4 & 29 & 4 \\
001201 & 东瑞股份 & 31.2 & 21.3 & -1.61 & 43.5 & -4.9 & 43 & 0 \\
300313 & 天山生物 & 24.4 & 2.0 & -0.27 & 20.7 & — & 28 & 2 \\
002234 & 民和股份 & 23.4 & 21.6 & -2.69 & 2.1 & -14.1 & 1 & 0 \\
600965 & *ST福成 & 22.0 & 11.6 & 0.57 & 5.4 & 2.6 & 0 & 0 \\
300967 & 晓鸣股份 & 20.2 & 12.2 & 0.85 & 21.2 & 10.8 & 33 & 0 \\
603717 & 天域生物 & 17.2 & 7.3 & -1.07 & 4.0 & -23.0 & 30 & 1 \\
002982 & 湘佳股份 & 16.2 & 44.7 & 0.33 & 7.2 & 1.9 & 15 & 0 \\
\end{longtable}
\endgroup


\subsection{生猪养殖：成本曲线决定生存}
\label{sec:breeding-hog}

\num{11} 家生猪养殖公司的 2025 年业绩呈现出教科书式的\hl{“头部盈利、尾部巨亏”}结构：

\begin{itemize}[leftmargin=1.4em, itemsep=2pt]
  \item \textbf{牧原股份}：营收 \num{1441.4} 亿元（两年 CAGR \num{14.0}\%）、
        归母净利润 \num{154.9} 亿元、ROE \num{20.6}\%，
        是全部 \num{104} 家农业上市公司中唯一净利润超过百亿元的企业；
  \item \textbf{温氏股份}：营收 \num{1038.6} 亿元、净利润 \num{52.7} 亿元、ROE \num{12.6}\%；
  \item \textbf{神农集团}：营收 \num{53.5} 亿元、净利润 \num{3.4} 亿元、ROE \num{6.9}\%，
        是中等规模企业中少见的正回报者；
  \item 尾部三家：\textbf{天邦食品}（\num{-13.1} 亿元）、\textbf{新五丰}（\num{-8.8} 亿元）、
        \textbf{正邦科技}（\num{-5.5} 亿元）合计亏损 \num{27.4} 亿元。
\end{itemize}

\hl{牧原与温氏合计实现 \num{207.5} 亿元净利润，而其余 \num{20} 家养殖公司合计净亏损
\num{4.2} 亿元}——板块 \num{203.4} 亿元的利润总额，\hl{全部来自这两家公司}。
这与第 \ref{sec:hog-cost} 节的判断一致：在猪价 \num{10}---\num{12} 元/公斤的深度亏损区间，
只有完全成本低于 \num{13} 元/公斤的企业还能盈利。
牧原的 \hl{自繁自养一体化与低豆粕日粮}结构使其成本处于行业最低分位，
而“公司 + 农户”模式下外购仔猪比例更高的企业，在仔猪价格从 \num{108.6} 元/公斤
跌至 \num{22.05} 元/公斤的过程中承受了双向亏损（仔猪跌价损失 + 育肥亏损）。

\hl{正邦科技是观察周期破坏力的最佳样本}：流通市值仍高达 \num{227} 亿元（板块第三），
但 2025 年营收仅 \num{147.9} 亿元、亏损 \num{5.5} 亿元、ROE \num{-4.9}\%。
市场对其给予高市值的原因是重整预期与产能保留价值，
而不是当前的盈利能力——\emph{这本身是周期底部的典型定价特征}。

\subsection{肉鸡养殖：反转最快的一环}
\label{sec:breeding-poultry}

\num{8} 家肉鸡养殖公司中，\textbf{圣农发展}（\num{13.80} 亿元、ROE \num{12.6}\%）
与 \textbf{立华股份}（\num{5.78} 亿元、ROE \num{6.3}\%）实现盈利，
而 \textbf{益客食品}（\num{-2.80} 亿元）、\textbf{民和股份}（\num{-2.69} 亿元）亏损。
这一分化与第 \ref{sec:other-poultry} 节的产业数据相互印证：
2025 年白羽肉鸡祖代更新量创历史新高、出栏均价跌至近十年低位，
\hl{上游种鸡与商品代养殖环节同时承压}；
而圣农发展以全产业链（自有种源 + 自养 + 屠宰 + 深加工）模式
锁定了下游食品端的利润，抵消了商品代的周期波动。

2026 年 1---4 月毛鸡养殖利润同比 +\num{297.22}\%，
\hl{禽链的盈利反转已经开始}，且由于产能周期短，
效益改善向财务报表的传导速度将快于生猪养殖。

\begin{keybox}[养殖业的三点判断]
\normalsize
\textbf{① 集中度提升不可逆。}牧原与温氏合计利润为正、其余公司合计为负，
行业资本回报向低成本龙头集中的趋势在 \num{2023}---\num{2025} 年持续强化。\par\vspace{3pt}
\textbf{② 高负债 + 亏损 = 融资依赖。}板块内 \num{6} 家资产负债率超过 \num{65}\%，
其中多家同时处于亏损状态，\hl{其对股权融资的依赖程度高于其他子行业}
（详见第 \ref{sec:financing} 节）。\par\vspace{3pt}
\textbf{③ 2026 年是养殖板块的“利润修复年”，但不是“估值修复年”}：
禽链反转快、猪链磨底慢，两者的报表差异会在 2026 年年报中明显拉开。
\end{keybox}
